\section{Artifact}
\label{app:artifact}

The artifact is one repository, to be released under Apache-2.0 at tag
\texttt{paper-v1}. It contains:
\begin{itemize}
\item the specification: \texttt{spec/layer.md} (normative), the A2A and
  libp2p-discovery bindings, the payment-rail mapping, and a changelog
  from version~1;
\item the Rust implementation (\texttt{map-rs}) and the Go implementation
  (\texttt{map-go}), the test vectors and their generator, and the
  28-cell interop matrix (\texttt{interop/run-matrix.ps1});
\item the simulator (\texttt{map-rs/sim}), the pre-registration
  (\texttt{eval/HYPOTHESES.md}, tag \texttt{prereg-v1}, amendments
  appended), the configurations, and the analysis;
\item every per-seed CSV of the runs reported here, with toolchain and
  CPU metadata (\texttt{eval/results/5bc7702/}).
\end{itemize}
\texttt{eval/run\_all.ps1} reproduces every table and figure. It first
checks golden hashes and that fast and real signing agree. On one 16-thread
desktop the full run takes about four hours, most of it the W2 trust grid.
The paper's figures come from \texttt{eval/analysis/make\_figs.py} and
\texttt{paper/figs/make\_flow.py}.

\section{Amendments to the Pre-registration}
\label{app:amendments}

All amendments are appended, dated and committed. The pre-registered text
above them is unchanged since \texttt{prereg-v1}.
\begin{description}
\item[A1] E1 uses fast signing, which the plan had reserved for larger
  runs. A check grid covering all eight arms shows a byte-identical CSV
  in real and fast modes, and every run repeats the check.
\item[A2] Clarification: every E1 arm gets the same fallback of up to
  three attempts, so arms differ only in ranking and filtering.
\item[A3] Clarification: B1 keeps decay and the per-issuer cap. This is
  a stronger baseline than an unweighted sum.
\item[A4] Repair of the D5 tie clause. As written, it was satisfied almost
  always and would have chosen W0. The repair was made after two-seed
  plumbing numbers had been seen; both readings are reported.
\item[A5] The grinding rate is set to the benchmark's 60,901 attempts per
  second on one core. A 16-core variant is run as an exploratory
  addition.
\item[A6] D5 result: W2 wins at $\alpha{=}0.2$. W1/W2 were re-implemented
  deterministically; the test compares them with the direct definitions
  for exact equality.
\item[A7] Implementation only: cache-blocked W2, bit-identical, after an
  interrupted run.
\item[A8] Exploratory: W2 collapses under camouflage at $\alpha{=}0.5$
  (0.551), while W1 does not (0.206).
\item[A9] Minimax re-test of the recommended weighting, pushed before its
  run. W1 wins with a worst case of 0.200; W0 reaches 0.348 and W2 0.550.
\item[A10] Exploratory, pushed before its run: the number of Sybils per
  adversarial provider, swept at $\alpha{=}0.2$. W1 loses to W0 from 80
  Sybils per provider (6.4 per honest requester), and W2 becomes worse
  than random at 320. It adds a Sybil-ratio caveat to the W1
  recommendation.
\end{description}

\section{Recommended Parameters}
\label{app:params}
\begin{itemize}
\item \textbf{Trust:} decay 0.98/day, prior $(1,1)$, future-skew bound
  120\,s, per-issuer cap 50, issuer weighting W1.
\item \textbf{Checkpoints:} interval $c{=}8$ Parts.
\item \textbf{Discovery:} $r{=}3$ sampled shards; D-6 filter mandatory;
  $d{=}3$ disjoint lookup paths; status pulled on demand; descriptor
  validity of at most 24\,h.
\end{itemize}
The trust values were fixed before the runs and were not tuned. $c{=}8$ and W1 were chosen from the results, as described in \S\ref{sec:eval-checkpoint} and \S\ref{sec:eval-trust}.
